\section{Ordered exponential sums and ordinary root estimates}
\label{ordsumrepair:section}

The proof of \cite[Lemma~5.3, pp.~31--33]{KoukLocal} uses a real
threshold followed by integer bands. We give the proof with an explicit
integer-rounding correction in this band decomposition and recover the
$A=0$ case of \cite[Lemma~5.5]{Kouk} used below. The two order-statistics
inputs are stated next. All constants below may depend on fixed $C>1$
and $\mu>1$.

Put
\begin{align*}
 \Delta_r&=\{0\leq\xi_1\leq\cdots\leq\xi_r\leq1\},\\
 S_r(u,v)&=\{\xi\in\Delta_r:\xi_j\geq(j-u)/v\ (1\leq j\leq r)\},\\
 Q_r(u,v)&=r!\operatorname{vol}(S_r(u,v)).
\end{align*}

\begin{lemma}[Flat-wall upper bound]
\label{ordsumrepair:flat-input}
For an integer $r\geq1$ and real $u,v>0$ with $w=u+v-r>0$,
\[
 Q_r(u,v)\ll\frac{(u+1)(w+1)}r.
\]
\end{lemma}

This is the upper assertion of \cite[Lemma~5.1, p.~30]{KoukLocal}.
Koukoulopoulos derives this statement from Ford's estimates
\cite[Theorem~1]{FordSmirnov2008} and
\cite[Lemma~11.1]{FordDivisors2008}.
The lower assertion of Lemma~5.1 is not used here.

\begin{lemma}[Crowding estimate]
\label{ordsumrepair:crowding-input}
Fix $C>1$. Suppose $g,r,s,u,v$ are integers satisfying
\[
 2\leq g\leq r/2,\quad s\geq0,\quad r\leq Cv,\quad
 u\geq0,\quad u+v\geq r+1.
\]
Let $R$ consist of the points of $S_r(u,v)$ for which there is an
integer $g+1\leq l\leq r$ such that
\[
 \frac{l-u}{v}\leq\xi_l\leq\frac{l-u+1}{v},\qquad
 \xi_{l-g}\geq\frac{l-u-s}{v}.
\]
Then
\[
 \operatorname{vol}(R)\ll_C
 \frac{[C(s+1)]^g}{(g-1)!}
 \frac{(u+1)(u+v-r)}{(r+1)!}.
\]
\end{lemma}

This is \cite[Lemma~5.2, p.~30]{KoukLocal}. That source obtains it
from \cite[Lemma~4.3]{FordDyadic2008}, using Lemma~5.1 in place of
Ford's Lemma~4.1.
The stated changes are the fixed constant $10$ to $C$, the factor
$(u+v-r)^2$ to $u+v-r$, and $(g-2)!$ to $(g-1)!$.
We use these two order-statistics estimates in precisely the forms
stated above.

For any $\gamma\geq0$, define
\[
 \mathcal T_\mu(r,v,\gamma)=
 \left\{\xi\in\Delta_r:
    \sum_{i\leq j}\mu^{v\xi_i}\geq\mu^{j-\gamma}
                           \ (1\leq j\leq r)\right\}.
\]

\begin{lemma}[Corrected integer-parameter ordered-sum estimate]
\label{ordsumrepair:integer-bound}
Fix $C>1$ and $\mu>1$. For integers $r,v,\gamma$ satisfying
$\gamma\geq0$ and $1\leq r\leq Cv$, set $b=r-v$, $d=b-\gamma$, and
\[
 Y=\begin{cases}b,&d\geq1,\\
                 (1-d)(\gamma+1),&d\leq0.
   \end{cases}
\]
Then
\begin{equation}
 \operatorname{vol}(\mathcal T_\mu(r,v,\gamma))
 \ll_{C,\mu}\frac{Y}{\mu^{\mu^d}(r+1)!}.
 \label{ordsumrepair:integer-estimate}
\end{equation}
\end{lemma}

\begin{proof}
The only change of threshold from the source proof is
\begin{equation}
 B=\left\lceil2+\frac{\log32}{\log\mu}\right\rceil,
 \qquad t=\max\{d,B\}.
 \label{ordsumrepair:threshold}
\end{equation}
Both $d$ and $t$ are integers. For $\xi\in\mathcal T_\mu(r,v,\gamma)$
let $D=\max_{j\leq r}(j-\gamma-v\xi_j)$.
If $D<t$, all coordinates satisfy
$\xi_j>(j-\gamma-t)/v$; call this part $V_1$.
Otherwise choose the integer $h=\lfloor D\rfloor+1\geq t+1$ and
an index $l$ attaining $D$. Then
\begin{equation}
 \min_{j\leq r}\left(\xi_j-\frac{j-\gamma}{v}\right)
 =\xi_l-\frac{l-\gamma}{v}
 \in[-h/v,(-h+1)/v].
 \label{ordsumrepair:band}
\end{equation}
Call the union of these bands $V_2$. This covers also $D=t$;
any overlaps at endpoints are harmless. The unrounded threshold
$t_0=\max\{d,2+\log32/\log\mu\}$ would omit points with
$t_0<D<\lceil t_0\rceil$ when $t_0$ is not an integer.

If $d\geq B$, then $t=d$ and $V_1$ is empty: the inequality for
$j=r$ would give $\xi_r>1$.
If $d<B$, then $t=B$, and Lemma~\ref{ordsumrepair:flat-input}
applies with $u=\gamma+B>0$, $w=B-d\geq1$. Thus
\begin{equation}
 \operatorname{vol}(V_1)
 \ll\frac{(\gamma+B+1)(B-d+1)}{(r+1)!}.
 \label{ordsumrepair:v1}
\end{equation}
Here $1/(r r!)\leq2/(r+1)!$. For $d\leq0$ the numerator is at most
$(B+1)^2Y$ and $\mu^{\mu^d}\leq\mu$. For $1\leq d<B$ it is at
most $B(B+2)b$, while $\mu^{\mu^d}\leq\mu^{\mu^{B-1}}$.
Since $B$ depends only on $\mu$, both cases give the required
bound for $V_1$, including its double exponential denominator.

Fix an integer band $h\geq t+1$. Write
\[
 a=\left\lceil\frac{\log4}{\log\mu}\right\rceil,\qquad
 m_0=1+a,\qquad h_0=h-m_0.
\]
These are integers, and
\begin{equation}
 h_0\geq t-\left(\frac{\log4}{\log\mu}+1\right)
       \geq1+\frac{\log8}{\log\mu}.
 \label{ordsumrepair:h0}
\end{equation}
We claim that there is an integer $m\geq h_0$ with
$g_m=\lfloor\mu^m\rfloor<l/2$ such that
\begin{equation}
 \xi_{l-g_m}\geq\frac{l-\gamma-2m}{v}.
 \label{ordsumrepair:cluster}
\end{equation}
For completeness, order the last half of the first $l$ summands and
group their backward indices into $[g_m,g_{m+1})$. Enlarging the first
group to $\mu^{h_0}$ terms and every subsequent group to
$g_{m+1}-g_m$ terms gives
\[
 \sum_{j\leq l}\mu^{v\xi_j}
 \leq2\left(\mu^{h_0}\mu^{v\xi_l}
       +\sum_{\substack{m\geq h_0\\g_m<l/2}}
            (g_{m+1}-g_m)\mu^{v\xi_{l-g_m}}\right).
\]
If \eqref{ordsumrepair:cluster} failed for all candidates, the
band condition would therefore give
\[
 \sum_{j\leq l}\mu^{v\xi_j}
 <2\mu^{l-\gamma}\left(\mu^{-a}
           +\sum_{m\geq h_0}(g_{m+1}-g_m)\mu^{-2m}\right)
 \leq\mu^{l-\gamma}.
\]
Indeed, $\mu^{-a}\leq1/4$, and summation by parts with
$\mu^m-1\leq g_m\leq\mu^m$ gives
\[
 \sum_{m\geq h_0}(g_{m+1}-g_m)\mu^{-2m}
 \leq\mu^{1-h_0}+\mu^{-2h_0}
 \leq2\mu^{1-h_0}\leq1/4.
\]
When there are no candidates, the first group alone bounds the sum
by $\mu^{l-\gamma}/2$, so no strict-inequality argument is needed.
Either case contradicts membership in $\mathcal T_\mu$ and proves
\eqref{ordsumrepair:cluster}.

Apply Lemma~\ref{ordsumrepair:crowding-input} with
$u=\gamma+h$, $s=2m$, $g=g_m$. The minimal-band condition places
$\xi$ in $S_r(u,v)$ and gives the required interval for $\xi_l$;
moreover
\[
 \xi_{l-g}\geq(l-\gamma-2m)/v\geq(l-u-s)/v.
\]
All parameters are integers, $g\geq2$ by
\eqref{ordsumrepair:h0}, $g<l/2\leq r/2$, $l\geq g+1$,
$s,u\geq0$, $u+v-r=h-d\geq1$, and $r\leq Cv$.
The external lemma already takes the union over $l$, so there is no
additional factor $r$. Summing only over $h,m$ yields
\begin{equation}
 \operatorname{vol}(V_2)\ll_C\frac1{(r+1)!}
 \sum_{h\geq t+1}\sum_{m\geq h-m_0}
 (\gamma+h+1)(h-d)
 \frac{[C(2m+1)]^{\lfloor\mu^m\rfloor}}
      {(\lfloor\mu^m\rfloor-1)!}.
 \label{ordsumrepair:v2sum}
\end{equation}
For fixed $C,\mu$, Stirling's lower bound implies
\[
 \frac{[C(2m+1)]^{\lfloor\mu^m\rfloor}}
      {(\lfloor\mu^m\rfloor-1)!}
 \ll_{C,\mu}\mu^{-\mu^{m+m_0}}.
\]
Its logarithm is
$-(\log\mu)m\mu^m+O_{C,\mu}(\mu^m\log(m+1)+m)$, which is eventually
smaller than $-(\log\mu)\mu^{m+m_0}$; the finitely many remaining
$m$ are absorbed into a constant depending only on $C,\mu$.
The successive ratios of $\mu^{-\mu^q}$ for $q\geq t+1\geq B+1$
are bounded above by a fixed number less than one. Thus the $m$ tail
is $\ll_\mu\mu^{-\mu^h}$. Writing $h=t+1+j$, each of the two
linear factors grows by at most a factor $j+1$, and the same
geometric domination sums $(j+1)^2$. Consequently
\begin{equation}
 \operatorname{vol}(V_2)\ll_{C,\mu}
 \frac{(\gamma+t+2)(t+1-d)}{\mu^{\mu^{t+1}}(r+1)!}.
 \label{ordsumrepair:v2}
\end{equation}
If $d\geq B$, its numerator is $b+2\ll b=Y$.
If $d\leq0$, it is at most $(B+2)(B+1)Y$.
If $1\leq d<B$, it is at most $B(B+3)b$.
In all cases $t+1\geq d$, which proves the desired bound for $V_2$.
Together with \eqref{ordsumrepair:v1} and the covering, this proves
\eqref{ordsumrepair:integer-estimate}.
\end{proof}

\begin{corollary}[The $A=0$ estimate for every count]
\label{ordsumrepair:a0}
Fix $\mu>1$. For all positive integers $r,v$ and real
$\gamma\geq\max\{0,r-v\}$,
\begin{equation}
 r!\operatorname{vol}(\mathcal T_\mu(r,v,\gamma))
 \leq C_\mu\min\left\{1,
      \frac{(\gamma-r+v+1)(\gamma+1)}r\right\}.
 \label{ordsumrepair:a0-bound}
\end{equation}
\end{corollary}

\begin{proof}
When $1\leq r\leq2v$, put $b=r-v$ and
$\Gamma=\lceil\gamma\rceil$. The event increases with $\gamma$, so
$\mathcal T_\mu(r,v,\gamma)\subseteq\mathcal T_\mu(r,v,\Gamma)$.
Lemma~\ref{ordsumrepair:integer-bound} applies with $C=2$ and
$d=b-\Gamma\leq0$. Discarding its denominator
$\mu^{\mu^d}\geq1$, and using
\[
 (\Gamma-b+1)(\Gamma+1)
 \leq4(\gamma-b+1)(\gamma+1),
\]
gives \eqref{ordsumrepair:a0-bound}, after also using the simplex
bound $\operatorname{vol}(\mathcal T_\mu)\leq1/r!$.
If $r>2v$, then $b>r/2$, $\gamma\geq b$, and
$(\gamma-b+1)(\gamma+1)/r>1/2$. The simplex bound alone proves
the assertion with constant two in this case.
Thus no restriction on $r/v$ or centrality of the count is imposed.
The singular expression $1/r$ is never used at $r=0$.
\end{proof}

\subsection{The tail of the ordinary rectangle}
\label{ordsum:section}

The ordinary root rectangle involves a sum of exponentials of ordered
positions, rather than a flat order-statistics wall. We state precisely
the locally repaired input and its application to that rectangle.

\begin{lemma}[Ordered-sum estimate in the required $A=0$ range]
\label{ordsum:published}
Fix $\mu>1$. For positive integers $n,v$ and
$\gamma\geq0$ with $\gamma\geq n-v$, let
\[
 \mathcal T_\mu(n,v,\gamma)=
 \left\{0\leq\xi_1\leq\cdots\leq\xi_n\leq1:
      \sum_{\ell\leq j}\mu^{v\xi_\ell}
                   \geq\mu^{j-\gamma}\ (1\leq j\leq n)\right\}.
\]
Then
\begin{equation}
 n!\operatorname{vol}(\mathcal T_\mu(n,v,\gamma))
 \leq C_{\mu}\min\left\{1,
       \frac{(\gamma-n+v+1)(\gamma+1)}n\right\}.
 \label{ordsum:published-bound}
\end{equation}
\end{lemma}

\begin{proof}
This is Corollary~\ref{ordsumrepair:a0}, with $r=n$.
\end{proof}

This is exactly the $A=0$ portion of \cite[Lemma~5.5]{Kouk}, with the theorem numbering of arXiv:1102.3236v3.
The proof in
Section~\ref{ordsumrepair:section} includes integer rounding of the
threshold and the passage to real $\gamma$. Its hypotheses
include every $n/v$; there is no central-count restriction. The factor
$n!$ converts ordered-simplex volume to the probability for the order
statistics of $n$ independent uniform points. No assertion for $A>0$
is needed here.

\begin{lemma}[All-count tail of the ordinary rectangle]
\label{ordsum:rectangle-tail}
Fix an integer $2\leq h\leq H$ and $C_0\geq1$. Set
$c=h-1$, $w=\log h$ and $\mu=h^{1/c}$. For $v\geq1$ an integer,
let $\xi_1\leq\cdots\leq\xi_n$ be the ordered independent uniform
points on $[0,1]$, and put
\begin{align}
 Z&=(n-v)\log h,\notag\\
 F_n&=\min_{0\leq j\leq n}
       h^{-j}\left(1+\sum_{\ell\leq j}\mu^{v\xi_\ell}\right)^c,
 & W_n&=\min\{1,C_0e^Z F_n\}.
 \label{ordsum:rectangle}
\end{align}
There exist constants $C,C_1$ depending only on $H,C_0$ such that,
for every $n\geq0$ and $\ell\geq0$,
\begin{equation}
 \mathbb P\{W_n>e^{-\ell}\mid n\}
 \leq\min\left\{1,
       \frac{C(1+\ell)(C_1+Z+\ell)_+}{1+n}\right\}.
 \label{ordsum:rectangle-bound}
\end{equation}
For $\ell<0$ the event is empty.
\end{lemma}

\begin{proof}
Suppose first $n\geq1$ and the event is nonempty. The term $j=0$
gives $F_n\leq1$, hence $Z+\ell+\log C_0>0$. Define
\[
 \gamma_0=\frac{Z+\ell+\log C_0}{\log h},
 \qquad \gamma=\gamma_0+\frac{\log2}{\log\mu}.
\]
For each $1\leq j\leq n$, the event in question implies
\[
 1+\sum_{a\leq j}\mu^{v\xi_a}>\mu^{j-\gamma_0}.
\]
The sum without the initial $1$ is at least $1$. Consequently
\[
 \sum_{a\leq j}\mu^{v\xi_a}
       >\tfrac12\mu^{j-\gamma_0}=\mu^{j-\gamma}.
\]
Here $\gamma\geq0$, and the exact identity $Z=(n-v)\log h$ gives
\[
 \gamma-n+v
   =\frac{\ell+\log C_0}{\log h}
      +\frac{\log2}{\log\mu}\geq0.
\]
Lemma~\ref{ordsum:published} therefore applies in its proved $A=0$ range.
Furthermore
\[
 \gamma-n+v+1\leq C(1+\ell),\qquad
 \gamma+1\leq C(C_1+Z+\ell)_+
\]
on the nonempty event, after choosing $C_1$ large in terms of $H,C_0$.
There are only finitely many $\mu=h^{1/(h-1)}$, $2\leq h\leq H$;
the constants from Corollary~\ref{ordsumrepair:a0} may therefore be
replaced by their maximum over this fixed finite set.
Its probability is at most the right side of
\eqref{ordsum:rectangle-bound}, since $1+n\leq2n$.

If $n=0$, then $F_0=1$ and $W_0=\min(1,C_0e^Z)$ is deterministic.
On its nonempty event $Z+\ell> -\log C_0$. Increasing $C_1$ once
makes the displayed right side equal to one. This proves the assertion
also at count zero. Finally $W_n\leq1$ makes the event empty for
$\ell<0$.
\end{proof}

In the actual root application $v$ is the number of completed harmonic
cells, of width $\eta=\log h/(h-1)$, and their budget is $U=\eta v$.
Uniform latent points in the cells give exactly the variables used in
\eqref{ordsum:rectangle}; rounding changes the geometric constant $C_0$
by a fixed factor. Thus $Z=n\log h-(h-1)U$ exactly. A bounded change
of physical budget is kept in the external root translate, as in
Section~\ref{positive:section}. If there are no cells, the completed
Poisson count is identically zero and is covered directly. The subsequent
Poisson maximum-atom and layer integrals use the same count.
